\documentclass{ltjsbook}
\usepackage{docmute} 
\usepackage{../../myStyle} 

\begin{document}
\chapter{名義尺度水準の尺度モデル}
\label{H09_dual}

多次元尺度構成法は距離行列から始め，順序尺度水準しか持たない情報であっても可視化など次元解析をすることが可能な技術でした。
では順序尺度よりも低い水準である，名義尺度水準の次元を見るような分析方法はあるのでしょうか。ここではそれをみていくことにしましょう。

\section{カテゴリに数値を与える}

名義尺度水準は，数字の付け方が名義的です。たとえば男性に1,女性に2という数字を割り振って集計する，というような数値の与え方が名義尺度水準です。
このとき，男性に2，女性に1であっても構いませんし，男性に33465，女性に8768432を割り当てても問題ありません。ここでの数字は数字としての意味があるのではなく，<33465>という文字列が男性というカテゴリに対応している，という一対一対応だけがわかっていればそれで良いのです。数字にラベルとしての意味しかありませんから，これを使って計算することはありません。

間隔尺度水準の量的なデータであれば，線形モデルが色々開発されて利用されてきました。因子分析も回帰分析も主成分分析も，あるいはそれらを統合した構造方程式モデリングも，全て線形モデルの世界です。
分散共分散行列を標準化した行列は相関行列といいますが，相関係数が$1.0$(あるいは$-1.0$)の状態を考えれば明らかなように，分散共分散行列(や相関行列)で表現されるのは変数の直線的な関係性に限った話なのです。

しかし心理学の場合は中庸が良いようなシーンも少なくありません。たとえば血圧と健康度の関係でいえば，低血圧でも高血圧でも不健康であり，ちょうどいい血圧が一番健康的だというのはすぐにわかります。このように中庸が良い場合は，散布図がU字型に現れてきます。散布図がU字型であれば，相関係数としては0近くになりますが，血圧と健康が無関係だとは誰もいえないでしょう。

これについて，ごくシンプルかつ具体的な例をあげてみましょう。血圧と頭痛の頻度について調査したとします。血圧は高い，普通，低いの3段階。頭痛の頻度は「ない」，「たまに」，「ときどき」，「いつも」の4段階です。表\ref{tbl::14_01}のような結果を見ると，この集計表に直線的な関係は確かになさそうですね。でも関係ないわけではない，と思いませんか。
\begin{table}[h]
    \begin{center}
        \caption{血圧と頭痛}
        \begin{tabular}{c|cccc}
            血圧 & ない & たまに & ときどき & いつも \\
            \hline
            高い & 0  & 0   & 3    & 2   \\
            普通 & 5  & 0   & 0    & 0   \\
            低い & 0  & 2   & 3    & 1   \\
        \end{tabular}
        \label{tbl::14_01}
    \end{center}
\end{table}

このような場合は，「血圧が高い人か低い人は，頭痛がある」という傾向が見て取れるはずです。しかしこのデータ，機械的に血圧が高いを1，普通を2，低いを3とし，同様に頭痛もない〜いつもを1,2,3,4として相関係数を計算すると，$r=0.165$になります。相関関係では傾向をうまく読み取れていないのです。

その根本的な原因は，もちろんスコアリングの方法にあります。共分散を計算するには間隔尺度水準以上の情報が必要で，今回のような3,4件法では相関係数を計算して良い数字ではありません。
ではポリコリック相関係数のように順序尺度水準の相関係数なら良いか，といいたいところですが，それでもまだ不十分です。なぜなら，血圧が高い方から低い方まで，順番が保存されてしまっているからです。

関係を導き出すには抜本的な改革が必要です。たとえば表\ref{tbl::14_02}のようにすればどうでしょうか。こうすると左下から右上にかけての直線的な関係がまだ見えてきます。
これは血圧の順番を「高い」「低い」「普通」に並べ替えたものです。
また，頭痛の順番も「いつも」と「ときどき」をひっくり返しました。
\begin{table}[h]
    \begin{center}
        \caption{並び替えられた「血圧と頭痛」}
        \begin{tabular}{c|cccc}
            血圧 & ない & たまに & いつも & ときどき \\
            \hline
            高い & 0  & 0   & 2   & 3    \\
            低い & 0  & 2   & 1   & 3    \\
            普通 & 5  & 0   & 0   & 0    \\
        \end{tabular}
        \label{tbl::14_02}
    \end{center}
\end{table}
順番を並び替えてしまったというのは，尺度水準的には数字を無視して「高い」「低い」「普通」というカテゴリーとして扱ったということになります。つまり\keytermL{名義尺度水準}{めいぎしゃくどすいじゅん}レベルにまで落として考えたのです。

このように並べ替えて，血圧のスコアを高い$\to3$，低い$\to2$，普通$\to1$とし，また頭痛の頻度をいつも$\to3$，ときどき$\to4$と付け替えて相関係数を計算すると，今度は$r=0.810$になりました。こちらの方が線形性は高いことが数字でも確認できました。

ここで行ったのは，すべての反応カテゴリをただの言葉だと考えて，付与された数字を無視して並べ替えたことになります。そうすることで，左下から右上にかけて線形性を高めることができました。しかし「高い」と「低い」，「低い」と「普通」の配置も等間隔である必要はなく，もっと直線性がはっきりするように配置してやってもよいのでは，というアイデアが浮かびます。

ここで考え方が反転していることに注意してください。一般的なリッカート尺度では，反応カテゴリに(シグマ法などで)数字をつけて，変数間の線形性を算出して意味を考えるのでした。
ここでは逆に，変数間の線形性を最大にするように反応カテゴリに数字をつけてやろう，という考え方です。データの直線性というのはデータの特徴を最も強調し解釈しやすい形です。そのようにデータを整えるためには，反応カテゴリにどういう数字を付与すれば良いか，と考えるのです。

分析者がデータを解釈しやすくする目的で，カテゴリに数値を与えること。これが名義尺度水準の分析をするときの発想であり，林の「数量化理論」と呼ばれる一連の分析手法群です。
少し寄り道になりますが，林の数量化理論についてみておきましょう。

\subsection{林の数量化理論}
数量化の研究をしたのは，日本の偉大な統計学者，林知己夫(はやし ちきお)\footnote{林 知己夫(はやし ちきお，1918年6月7日 - 2002年8月6日)は，日本の統計学者。正四位勲二等，理学博士。統計数理研究所第7代所長。 社会調査・世論調査におけるサンプリング方法の確立を始め，数量化理論(Hayashi's Quantification Methods)の開発とその応用で知られる。1990年代以降，データの科学(Data Science)を提唱し，その研究・思想は現在へと引き継がれている。Wikipediaより。}という人です。その名を冠して林の数量化理論と呼ばれることがあります。

林はさまざまな研究成果を挙げており，論文ごとにデータに必要な分析方法を開発するというようなスタイルでした。その膨大な研究業績を，弟子である飽戸弘\footnote{飽戸 弘(あくと ひろし，1935年3月14日 - )は，日本の社会学者，東京大学名誉教授。専門は，社会心理学，コミュニケーション論。Wikipediaより。}が分類して，同じような分析方法ごとにI類，$\2$類，$\3$類，$\4$類，と呼んでいきました。

\begin{table}[h]
    \caption{林の数量化}
    \begin{tabular}{c|c|c|c|c}
        手法    & 外的基準 & データ     & 目的           & 関連する手法       \\
        \hline
        数量化I類 & 量的変数 & 質的変数    & 外的基準の予測      & 重回帰分析        \\
        数量化$\2$ 類 & 質的変数 & 質的変数    & 外的基準の判別      & 判別分析         \\
        数量化$\3$類 & なし   & 質的変数    & 変数間の関係の要約と記述 & 主成分分析・正準相関分析 \\
        数量化$\4$類 & なし   & 対象間の類似度 & 対象間の関係の要約と記述 & 多次元尺度法       \\
    \end{tabular}
    \label{tbl::14_03}
\end{table}

表\ref{tbl::14_03}にあるように，基本的には質的変数，すなわち名義尺度水準や\keytermL{順序尺度水準}{じゅんじょしゃくどすいじゅん}のデータが得られたときに，それを解釈するためにはどのような数値を割り振ってやれば良いか，という発想から生まれたものになっています。

さきほどの表\ref{tbl::14_01}のようなデータの並べ替えについては，数量化でいうところの$\3$類に該当します。外的基準を持たずに，変数間関係の要約と記述をするのですから，因子分析や主成分分析の名義尺度版だと考えてもいいかもしれません。ポイントは，データセットがカテゴリカル変数なので，数字の直線性を考えるのではなく，カテゴリの並べ替えも認めた上で尺度化するという，「分析者が主体的に数字を割り振る」という考え方が前面に打ち出されているところです。

カテゴリカル変数を対象にしたモデル群なので，分析の応用範囲は多岐にわたります。どのような変数でも名義尺度水準に落とすことができるからです。表\ref{tbl::14_04}には一般的なデータセットの形を示しました。IDがあって，Q1,Q2...と項目が列方向に並ぶ形です。一行が一人の反応を表しています。Q1がたとえば性別などの名義尺度水準の変数であっても，男性$\to1$，女性$\to2$のようにコード化するルールを決めて入力します。Q2はたとえばリッカート法で当てはまる，やや当てはまる・・・などの5件法だったとしましょう。その場合はシグマ法によるスコアを入れる，あるいは簡便的に当てはまるを5，やや当てはまるを4，といったように数字を割り振って入力しますね。

これをカテゴリカル変数と見なしてデータセットにした例が表\ref{tbl::14_05}です。IDは分析対象ではありませんから横に置くとして，Q1が2列に増えています。ID=1の人は男性なのですが，この人はQ1:Male=1かつQ1:Female=0というようにコード化されています。同様に，Q2には「どちらとも言えない」を選択しているのですが，これを3とするのではなく\texttt{00100}と5列にわたってコード化しているのです。このようにすることで，すべての変数をカテゴリーとして扱うことができるようになります。

\begin{table}[h]
    \begin{minipage}[t]{.30\textwidth}
        \begin{center}
            \caption{一般的なデータセット}
            \begin{tabular}{cccc}
                ID       & Q1       & Q2       & $\cdots$ \\\hline
                1        & 1        & 3        & $\cdots$ \\
                2        & 2        & 4        & $\cdots$ \\
                $\vdots$ & $\vdots$ & $\vdots$ &          \\\hline
            \end{tabular}
        \end{center}
        \label{tbl::14_04}
    \end{minipage}
    %
    \hfill
    %
    \begin{minipage}[t]{.60\textwidth}
        \begin{center}
            \caption{カテゴリカル化したデータセット}
            \begin{tabular}{ccccccc}
                ID       & Q1:M      & Q1:F     & Q2:5 & Q2:4 & Q2:3 & $\cdots$ \\\hline
                1        & 1         & 0        & 0    & 0    & 1    & $\cdots$ \\
                2        & 0         & 1        & 0    & 1    & 0    & $\cdots$ \\
                $\vdots$ & $\vdots $ & $\vdots$ &                               \\\hline
            \end{tabular}
        \end{center}
        \label{tbl::14_05}
    \end{minipage}
\end{table}

また表\ref{tbl::14_01}を並べ替えた時のように，こうした名義尺度のデータの線形性が最大になるように，行だけでなく，列も並べ替えます。データの中で線形性が最大になるように，反応カテゴリ\textbf{と}回答者に数字を割り振るのです。
ちなみに表\ref{tbl::14_01}は集計されたデータであり，今回の表\ref{tbl::14_05}は集計前のデータです。カテゴリカルな変数の分析の場合は，どちらでも良いのです。行と列の関係が表されている数字であれば，「ある/ない」のような二値反応でも，集計された度数でも構いません。さらにいえば，行と列の関係が記述されていればなんでもいいのです。

こうしてカテゴリカルな変数として表されたデータセットに対し，並べ替えも許して適切な数字を割り振るアルゴリズムのひとつに，交互平均法と呼ばれるものがあります。

\subsection{交互平均法}
具体的に，どのようなアルゴリズムで行・列に対する重みをつけるか見てみましょう。例えば，以下のような分割表が得られたとします（表\ref{DAexample}）。
\begin{table}[ht]
    \begin{center}
        \caption{データ例}
        \begin{tabular}{c|ccc|c}
            血圧 & 頭痛なし & たまに頭痛 & 頭痛あり & 計  \\
            \hline
            低い & 9    & 2     & 3    & 14 \\
            普通 & 3    & 1     & 4    & 8  \\
            高い & 6    & 2     & 5    & 13 \\
            \hline
            計  & 18   & 5     & 12   & 35 \\
        \end{tabular}
        \label{DAexample}
    \end{center}
\end{table}
ここで各カテゴリに値を割り振る手続きは，次のような手順で行います。
\begin{enumerate}
    \item 行方向の「なし($X_{1}$)」を1点，「たまに($X_{2}$)」を0点，「あり($X_{3}$)」を-1点と仮におく。
    \item 列方向の「低い($Y_{1}$)」，「普通($Y_{2}$)」，「高い($Y_{3}$)」を算出する。すなわち，
          $$Y_{1}=\frac{1 \times 9 + 0 \times 2 + (-1) \times 3 }{14} = 0.4285$$
          $$Y_{2}=\frac{1 \times 3 + 0 \times 1 + (-1) \times 4 }{8} = -0.125$$
          $$Y_{3}=\frac{1 \times 6 + 0 \times 2 + (-1) \times 5 }{13} = 0.0769$$
    \item Yの平均値を0にする。すなわち，
          $$ \bar{Y} = \frac{14 \times Y_{1} + 8 \times Y_{2} + 13 \times Y_{3}}{35}=0.1714$$
          $$Y_{1} = Y_{1} - \bar{Y} = 0.4285 - 0.1714 = 0.2571$$
          $$Y_{2} = Y_{2} - \bar{Y} = -0.125 - 0.1714 = -0.2964$$
          $$Y_{3} = Y_{3} - \bar{Y} = 0.0769 - 0.1714 = -0.0945$$
    \item Yの大きさを整える。すなわち，絶対値最大の要素で割って，全体が1を超えないようにする。
          $$Y_{1} = Y_{1} / maxY = 0.2571 / |-0.2964| = 0.8674 $$
          $$Y_{2} = Y_{2} / maxY = -0.2964 / |-0.2964| = -1.000 $$
          $$Y_{3} = Y_{3} / maxY = -0.0945 / |-0.2964| = -0.3188 $$
    \item これらの重みを使って，今度は行の重みを算出する。
          $$X_{1}=\frac{Y_{1} \times 9 + Y_{2} \times 3 + Y_{3} \times 6 }{18} = 0.1608$$
          $$X_{2}=\frac{Y_{1} \times 2 + Y_{2} \times 1 + Y_{3} \times 2 }{5} = 0.0194$$
          $$X_{3}=\frac{Y_{1} \times 3 + Y_{2} \times 4 + Y_{3} \times 5 }{12} = -0.2493$$
    \item Xの平均値を0にする。
          $$ \bar{X} = \frac{18 \times X_{1} + 5 \times X_{2} + 12 \times X_{3}}{35}=-3.806e-17$$
          $$X_{1} = X_{1} - \bar{X} = 0.1608$$
          $$X_{2} = X_{2} - \bar{X} = -0.0194$$
          $$X_{3} = X_{3} - \bar{X} = -0.2493$$
    \item Xの大きさを整える。
          $$X_{1} = X_{1} / maxX = 0.6450 $$
          $$X_{2} = X_{2} / maxX = 0.0780 $$
          $$X_{3} = X_{3} / maxX = -1.000 $$
    \item 2〜7を収束するまで繰り返す。収束判定は，
          $$ \delta_{Y} = \sum Y_{i}^{t+1} - Y_{i}^{t} $$
          $$ \delta_{X} = \sum X_{i}^{t+1} - X_{i}^{t} $$
          $$ として，\delta_{X} < \epsilon　且つ　\delta_{Y} < \epsilon $$
\end{enumerate}

このように，適当な数字から始めても交互に平均(項目の中心)を合わせていくことによって，最終的に次のような最適な重みが得られます。
\[
    X=
    \left(
    \begin{array}{c}
            0.644 \\
            0.081 \\
            -1.00 \\
        \end{array}
    \right)
    ,Y=
    \left(
    \begin{array}{c}
            0.851 \\
            -1.00 \\
            -0.301
        \end{array}
    \right)
\]

これが交互平均法による重み付けアルゴリズムです。

この方法をコツコツとあらゆるデータに応用していくことももちろんできますが，数学的にもう少し洗練された方法を使うことが一般的です。すなわち，行列演算を用いるアプローチです。

これまでの回帰分析，因子分析からSEMに至るまでの流れは，変数間関係だけを考えてきました。N行M列のデータ行列の計算の途中で，行に関する情報は平均化して潰されてしまい，最終的には$M \times M$サイズの正方行列だけを扱うことになったのでした。そしてM個の変数に\keytermL{因子負荷量}{いんしふかりょう}などの重みをつけて考察してきました。
これを数学的な面から説明すると，正方行列の場合は\keytermL{固有値分解}{こゆうちぶんかい}をして，\keytermL{固有値}{こゆうち}と\keytermL{固有ベクトル}{こゆうべくとる}を求め，固有ベクトルが変数の重みになるのでした。
固有値分解とは
\[ \bm{Ax} = \lambda\bm{x} \]
の関係を見つけることでした。一般に$n\times n$の行列からは$n$個の固有値・固有ベクトルのセットが得られ，実対称行列であれば固有ベクトルからなる行列と，固有値を対角に持つ行列を使って，
\[ \bm{A} = \bm{X \Lambda X}' \]
と書くことができます。ここで固有ベクトルからなる行列$\bm{X}$を縦・横に並べて掛け合わせたものが$\bm{A}$になるのは，$\bm{A}$が縦・横同じサイズの正方行列だったからです。

しかし今回のカテゴリカルなデータは，$N\times M$行の矩形行列になります。
矩形行列の場合は縦と横の長さが違いますから，同じベクトルを縦・横にしてかけたところでサイズが整いません。そこで
\[ \bm{B} = \bm{Y M X}' \]
のように分解します。ここで$\bm{B}$は$N \times M$行列とすると，$\bm{Y}$は$N$行，$\bm{X}$は$M$列の行列です。N行M列のデータセットですから，行のベクトルと列のベクトルそれぞれが計算されることになります。また$\bm{M}$は対角に$\rho_i$を含んだ正方行列であり，この値は\keyterm{特異値}{とくいち}{Singular balue}とよばれます。矩形行列版の固有値のようなものですね。

このように分解する方法は\keyterm{特異値分解}{とくいちぶんかい}{Singular value decomposition}と呼ばれ，行および列に対応する\keytermL{特異ベクトル}{とくいべくとる}をその重み，座標と考えることになります。ちなみに固有値分解は特異値分解の特殊な例(正方行列の場合)だと考えることができます。

\section{双対尺度法}
さて，この数量化$\3$類はおもしろいもので，同時期に独立にこの手法がフランス，カナダでも発展しており，2つの別名を持っています。フランス学派が開発した手法は\keyterm{対応分析}{たいおうぶんせき}{correspondence analysis}といい，カナダ在住の日本人，西里静彦\footnote{西里 静彦(にしさと しずひこ，1935年6月9日 - )は，カナダの行動計量学者（計量心理学）。学位はPh.D.（ノースカロライナ大学・1966年）。トロント大学名誉教授，アメリカ統計学会フェロー，日本行動計量学会名誉会員。北海道十勝郡浦幌町出身(札幌市生まれ)。Wikipediaより。}が開発した手法は\keyterm{双対尺度法}{そうついしゃくどほう}{Dual Scaling}といいます。

では数量化$\3$類とコレスポンデンス分析，および双対尺度法とよばれる一群の手法が，それぞれどのように異なっているのかについてみていきましょう。

\subsection{数量化$\3$類とコレスポンデンス分析}
数量化$\3$類とコレスポンデンス分析は，分析のもとになる分割表の考え方の違いに求めることができます。例えば，「はい」と「いいえ」で回答が得られる項目が3つあったとしましょう。このとき，数量化$\3$類が分析対象にする行列は，「はい」を1，「いいえ」を0とした分割表であり，表\ref{hayashi}のような表し方になります。
\begin{table}[ht]
    \begin{center}
        \caption{数量化が分析する行列}
        \begin{tabular}{c|ccc|c}
            被験者 & 項目1 & 項目2 & 項目3 & 合計 \\
            \hline
            S1  & 1   & 0   & 1   & 2  \\
            S2  & 0   & 0   & 1   & 1  \\
            S3  & 1   & 1   & 0   & 2  \\
            S4  & 1   & 0   & 1   & 2  \\
            S5  & 1   & 1   & 1   & 3  \\
            \hline
                & 4   & 2   & 4   &    \\
        \end{tabular}
        \label{hayashi}
    \end{center}
\end{table}

これに対して，コレスポンデンス分析と双対尺度法が分析対象とする行列は，表\ref{corepon}のように考えます。
\begin{table}[ht]
    \begin{center}
        \caption{コレスポンデンス分析，および双対尺度法が分析する行列}
        \begin{tabular}{c|cccccc|c}
            被験者 & 項目1(Yes) & 項目1(No) & 項目2(Yes) & 項目2(No) & 項目3(Yes) & 項目3(No) & 合計 \\
            \hline
            S1  & 1        & 0       & 0        & 1       & 1        & 0        & 3  \\
            S2  & 0        & 1       & 0        & 1       & 1        & 0        & 3  \\
            S3  & 1        & 0       & 1        & 0       & 0        & 1        & 3  \\
            S4  & 1        & 0       & 0        & 1       & 1        & 0        & 3  \\
            S5  & 1        & 0       & 1        & 0       & 1        & 0        & 3  \\
            \hline
                & 4        & 1       & 2        & 3       & 4        & 1        &    \\
        \end{tabular}
        \label{corepon}
    \end{center}
\end{table}

情報量としては同じですが，YesでなければNoであるとするか，Yesがある/Noがあるという形で示すかの違いでしかありません。この分析するもととなる行列の違いがあるだけで，それを特異値分解して行・列の両方に重みを与える点では同じ分析手法であるといえます。
\subsection{コレスポンデンス分析と双対尺度法}
コレスポンデンス分析と双対尺度法の違いは，もう少し複雑な説明になります。これらの違いについて考える前に，\keytermL{双対性}{そうついせい}について説明しなければなりません。

まず，双対尺度法は名義尺度水準同士の変数間関係を分解するものですから，表\ref{DAexample}のような分割表を分解することになります。
分割表で表される行と列がどれぐらい関係しあっているか，あるいは独立であるかを見るための指標として，$\chi^2$値があります。
双対尺度法は，この$\chi^2$値を次のように分解していくモデルでもあります。
\begin{equation}
    \chi^{2}=\chi_{1}^{2}+\chi_{2}^{2}+\chi_{3}^{2}+\ldots
    \label{chi}
\end{equation}

主成分分析が固有値分解したように，双対尺度法は特異値分解を使って分解していくのですが，これを要素レベルで見ていくと次のようになります。
まず分割表の$i$行$j$列目の度数を$f_{ij}$とし，行$i$の総和(周辺度数)を$f_{i.}$，列の総和を$f_{.j}$と表すことにします。また，総度数を$f_t$と書くとしますと，各行・各列が周辺度数の比率通りの独立な関係であるならば，
\[ f^{\ast}_{ij} = \frac{f_{i.}f_{.j}}{f_t} \]
となります。これは各セルの期待値を計算していることになります。$\chi^2$値は，これを各セルについて計算して，期待値と観測値のずれの総和，すなわち各セルの分散の大きさのようなものとして表していることになります。
\[ \chi^2 = \sum_{i=1}^m \sum_{j=1}^n \frac{(f_{ij}-f^{\ast}_{ij})^2}{f^{\ast}_{ij}} \]
双対尺度法は得位置分解を使って，行と列に最適な重みを計算するのですが，要素レベルで見ると次のような分解になります。
\[    
    f_{ij}=\underbrace{\underbrace{\underbrace{\frac{f_{i.}f_{.j}}{f_{t}}( 1}_{第0次近似}+\rho_{1}y_{1i}x_{1j}}_{第1次近似}+\rho_{2}y_{2i}x_{2j}}_{第2次近似}+\ldots+\rho_{k}y_{ki}x_{kj})
\]
第0次近似は期待値の成分であり，これは行と列が独立な場合の大きさですから分析に際しては意味がなく，無意味解といわれたりします。分析をして意味があるのは次の第1次近似からであり，このとき$y_{1i}$は$y$列の第1ウェイト，$x_{1j}$は$j$行の第1ウェイトといいます。また$\rho$を特異値といい，この二乗$\rho^2$が固有値とよばれます。特異値や固有値はその次元の重要度を表すというのは，因子分析の時と同じですね。

特にどの項までで近似するかを示すのに，第n次近似という表現を用いるあたりは因子分析とも同じです。またデータが$m \times n$の場合，第$k=min(m,n)-1$次近似でデータを完全に説明することができます。

さて双対の関係とは，行$Y$と列$X$の重みとして分解した式(\ref{soutui})にみられる，次の関係です。
\begin{equation}
    \label{soutui}
    \left \{
    \begin{array}{ccc}
        Y_{i} & = & \displaystyle{\frac{1}{\rho}\frac{\sum f_{ij}X_{j}}{f_{i.}}} \\
        \\
        X_{j} & = & \displaystyle{\frac{1}{\rho}\frac{\sum f_{ij}Y_{i}}{f_{.j}}} \\
    \end{array}
    \right.
\end{equation}

この双対の関係が意味するのは，$Y$のベクトルが張る空間と$X$のベクトルが張る空間は等しくなく，一方が他方に射影されるためには特異値をもちいて変換しなければならない，ということです。たとえばある回答者Y1が，X1,X2,X3という３つの項目を選択した場合，X空間に射影されるYの重みづけられた座標は，Xが作る三角形の重心に来ることになります。

双対尺度法をつかうと，行の空間を見出すことも，列の空間を見出すこともできますが，その両者は別の座標を表現しています。行変数と列変数を同時に表現できれば両者の関係がわかって良いのですが，行空間のなかに列変数をプロットするとか，列空間のなかに行変数をプロットすることを考えるときは，一方から他方への座標を変換しなければならないのです。双対尺度法の文脈では，行空間にプロットするのか列空間にプロットするのかを選択し，空間内での距離関係に注意して考えることが基本です。

ところがコレスポンデンス分析は，これをひとつの空間に射影します。$\rho Y_{i}$と$\rho X_{j}$は，大きさ(ノルム)を整えてあるので分散が同じですから，$\rho Y_{i}$と$\rho _{j}$をひとつの空間に射影しようじゃないか，というのがコレスポンデンス分析のやり方なのです。厳密には数学的に正しくないやり方かもしれませんが，両者共通の空間に同時にプロットできますから，行変数と列変数の対応が見やすいという利点があります。あくまでも，同じ空間にプロットできているわけではない点には，注意が必要です。

\section{双対尺度法の発展}
双対尺度法をつかうことで，名義尺度水準のデータを次元解析できるようになりました。
リッカート法でなんとなく1,2,3,4,5という数字をつけて，潜在的な次元を頂戴したとありがたがるのではなく，分析者が積極的に得られた反応パターンに数字を与える。その目的は得られたデータを最も説明しやすくする(線形性を最大にする)ことである，というのが「尺度化」「数量化」のいわんとすることです。

双対尺度法はデータや分析方法に工夫をすることで，他にも面白い分析をすることができます。いくつか双対尺度法の発展モデルを見てみましょう。
\subsection{順序データの場合}
例えば好きな政党を順に番号をつけてもらうといった，順序データがあったとしましょう（表\ref{jun}）。
この順序データを双対尺度法で分析するには，少し工夫するだけでOKです。
\begin{table}[ht]
    \begin{center}
        \caption{順序データの場合}
        \begin{tabular}{c|ccccc|c}
            被験者    & 自民党    & 民主党    & 公明党    & 自由党    & 社民党    & 合計     \\
            \hline
            S1     & 1      & 2      & 4      & 5      & 3      & 15     \\
            S2     & 2      & 1      & 3      & 4      & 5      & 15     \\
            \vdots & \vdots & \vdots & \vdots & \vdots & \vdots & \vdots \\
        \end{tabular}
        \label{jun}
    \end{center}
\end{table}

双対尺度法で分析するには，この順序関係が入った表を，\keyterm{ドミナンス表}{どみなんすひょう}{dominance table}に変換してから計算を行うことになります。ドミナンス表とは，$R_{ij}$を被験者$i$の項目$j$に対する順位とすると
\begin{equation}
    d_{ij}=n+1-2R_{ij}
\end{equation}
で表される数値が入った表のことです。これは，ある項目が他の残りの項目それぞれと比べて順位が上であれば+1，下であれば-1である数値で，いわば他の項目と何勝何敗であったかを表す数値になっています。双対尺度法では変数同士の相対的な結びつきの強さがロウデータになりますから，このように変換することで相対的な関係の強さに置き換えるわけです。今回の例では，ドミナンス表は以下のようになります(表\ref{dominance}）。
\begin{table}[ht]
    \begin{center}
        \caption{ドミナンス表}
        \begin{tabular}{c|ccccc|c}
            被験者    & 自民党    & 民主党    & 公明党    & 自由党    & 社民党    & 合計     \\
            \hline
            S1     & 4      & 2      & -2     & -4     & 0      & 0      \\
            S2     & 2      & 4      & 0      & -2     & -4     & 0      \\
            \vdots & \vdots & \vdots & \vdots & \vdots & \vdots & \vdots \\
        \end{tabular}
        \label{dominance}
    \end{center}
\end{table}

あとはこれを特異値分解して，行・列の両方に重みを与えることになります。このようなドミナンス数の考え方は，例えば一対比較法のデータにも応用できます。一対比較法とは，例えばA,B,Cの３つのカテゴリがあった場合に，AとBを比較してどちらが優位であったか，AとCは，BとCは，というように一対一での優劣(似ている，似ていないとか，好ましい，好ましくないなど比較の次元は問いません)を競いあったデータのことです。このようなデータであっても，双対尺度法を使えば変数同士の尺度化，回答者の同時プロットも可能にします。

くどいようですが，尺度構成の基本は一対比較です。何らかの次元で，一方が他方より大か小かを判断する次元があるとするならば，それを取り出そうとする考え方です。得られたデータの中から最大の説明をするように数値を割り振る双対尺度法は，新しい心理尺度の可能性を開いてくれるかもしれません。

\subsection{強制分類法}

少しトリッキーな技術ですが，\keyterm{強制分類法}{きょうせいぶんるいほう}{Forced Classification}という考え方も心理尺度の中では有用かもしれません。

例えば2つの項目，F1,F2があって，それがそれぞれ「はい」・「いいえ」で答えるデータであったとしましょう。データは行列$[\bm{F1},\bm{F2}]$のようになりますが，ここでF2の列を意図的に水増しすることを考えます。例えばF2を5つ並べて，
\[    [\bm{F1},\bm{F2},\bm{F2},\bm{F2},\bm{F2},\bm{F2}] = [\bm{F1},5\bm{F2}] \]
のようにすることもできます。これは，項目2を5回繰り返したもので，データとしては表\ref{FCexam}になっていることとおなじです。
\begin{table}[ht]
    \begin{center}
        \caption{繰り返された表}
        \begin{tabular}{cc|cc}
            x1 & x2 & x3 & x4 \\
            \hline
            1  & 0  & 5  & 0  \\
            1  & 0  & 0  & 5  \\
            0  & 1  & 0  & 5  \\
            0  & 1  & 5  & 0  \\
        \end{tabular}
        \label{FCexam}
    \end{center}
\end{table}

一般に，ある項目$j$を$k$回繰り返したとすると，それは項目$j$に重み$k$を掛け合わせたことと同じになります。では繰り返しを5回ではなく，10回，20回・・・・とどんどん増やしていくとしましょう。
そして増やした行列それぞれに対して双対尺度法を繰り返すとどうなるでしょうか。当然のことながら，今回の例ではF2についての反応，カテゴリでいうx3,x4の重要度が増していくことになりますから，構成される空間はF2が第一次元に強く強く関係していきます。
繰り返しが$\infty$に近づくにつれ，F2が双対尺度法で得られる第一次元の解に与える寄与率が1.00に近くなっていくわけです。これを言い換えれば，F2の軸を基準にして，他の項目がそれに合うように強制的に分類して行くことでもあります(図\ref{FC})。実際には無限の重みをつける必要はなく，十数倍の重みづけで十分です
\begin{figure}[ht]
    \begin{center}
        \begin{tikzpicture}
            \draw[very thick] (-5,0)--(5,0); %x軸
            \draw[-] (0,-5)--(0,5); %y軸
            \coordinate (X1) at (3,2);
            \coordinate (X2) at (-3,-2);
            \coordinate (X3) at (3,-2);
            \coordinate (X4) at (-3,2);
            \fill (X1) circle [radius=2pt];
            \fill (X2) circle [radius=2pt];
            \fill (X3) circle [radius=2pt];
            \fill (X4) circle [radius=2pt];
            \node[above right] at (X1) {$X_1$};
            \node[above] at (X2) {$X_2$};
            \node[below] at (X3) {$X_3$};
            \node[above] at (X4) {$X_4$};
            \coordinate (F1) at (1,3.5);
            \coordinate (F2) at (-1,-3);
            \coordinate (F3) at (4,0);
            \coordinate (F4) at (-4,0);
            \fill[red] (F1) circle [radius=2pt];
            \fill[red] (F2) circle [radius=2pt];
            \fill[blue] (F3) circle [radius=3pt];
            \fill[blue] (F4) circle [radius=3pt];
            \node[above] at (F1) {$X_1^{\infty}$};
            \node[above] at (F2) {$X_2^{\infty}$};
            \node[above right] at (F3) {$X_3^{\infty}$};
            \node[below] at (F4) {$X_4^{\infty}$};
            \draw[->,>=Latex,dashed,very thick,bend right = 30] (X1) to (F1);
            \draw[->,>=Latex,dashed,very thick,bend right = 30] (X2) to (F2);
            \draw[->,>=Latex,dashed,very thick,bend right = 30] (X3) to (F3);
            \draw[->,>=Latex,dashed,very thick,bend right = 30] (X4) to (F4);
        \end{tikzpicture}
        \caption{強制分類法による布置}
        \label{FC}
    \end{center}
\end{figure}

強制分類法は，何か明確な外的基準があるときに，そこを基準にして周りを見てみようという分析方法です。因子分析のように，目に見えない軸を取り出して「この軸は何だろう」と悩むのではなく，分析者の分析視点を明確にして「この観点から考えれば，対象はこのようにプロットされます」と論じる方が責任の所在がはっきりして良いかもしれません。応用例として，学術的な分類法や法的な基準があるような領域で，専門家と非専門家の評定がどの程度ずれているかといったことを考えるときに，専門家の判断パターンを基準に重みづけをし，非専門家がどのようなずれを生じさせるのかを考えるといった可能性が考えられます。

\subsection{全情報分析}
式\ref{soutui}に挙げた双対の関係にあるように，列の最適な重みづけによる平均値は行の重みに特異値をかけたもの，行の最適な重みづけによる平均値は列の重みに特異値をかけたもの，になります。つまり双対尺度法によるプロットは，「列空間」に行・列の各カテゴリをプロットするか，「列空間」にそれらをプロットするかのどちらかを選ばなければなりません。行の空間と列の空間は異なる縮尺を持つ空間なのです。

ここで行変数と列変数を同一の空間にプロットすることを考えたときに，行変数と列変数の空間的隔たりがどの程度あるかを次の式で算出することができます。
$$ d_{ij} = \sqrt{\sum_{k=1}^K \rho_k^2(y_{ik}^2+x_{jk}^2-2\rho_ky_{ik}x_{jk})} $$
また，行変数同士($y_i,y_{i*}$)の距離や，列変数同士の距離($x_j,x_{j*}$)は，
$$ d_{ii*} = \sqrt{\sum_{k=1}^K \rho_k^2(y_{ik}-y_{i*k})^2}$$
$$ d_{jj*} = \sqrt{\sum_{k=1}^K \rho_k^2(x_{jk}-x_{j*k})^2}$$
で求めることができます。これらを距離行列$\bm{D}_{yy} = \{ d_{ii*} \},\bm{D}_{xx} = \{ d_{jj*} \}, \bm{D}_{xy} = d_{ij}, \bm{D}_{yx} = d_{ji} $からなる行列としてまとめれば，次のようになります。
\begin{equation}
    \bm{D} =
    \begin{bmatrix}
        \bm{D}_{yy} & \bm{D}_{yx} \\
        \bm{D}_{xy} & \bm{D}_{xx}
    \end{bmatrix}
\end{equation}
これは行の空間，列の空間，行と列の空間(と列と行の空間)の全てが含まれた，大きなデータセットです。一方から他方の関係を考えるだけでなく，全ての組み合わせを含んだ行列です。この行列を分析することを\keyterm{全情報分析}{ぜんじょうほうぶんせき}{Total Infromation Analysis}といいます。この行列は距離行列ですから，クラスター分析などで類似性をみつける分類などが考えられます。


双対尺度法は，行にも列にも，データのあらゆる相にたいして重みをつけて分析することができます。もちろんカテゴリの数だけ，人の数だけ情報がでてくるので単純な話ではなくなります。各項目の各反応カテゴリに対して，最適な値が割り振られることになりはするのですが，その数が膨大になってしまうのです。例えば5件法で3変数分のデータを取って1因子が出た，というのであれば話は単純ですが，双対尺度法は15カテゴリに重みがつけられ，しかもそれらが必ずしも直線上に並んでいるわけでもなく，全体を眺めて意味を考えなければならないのです。しかも同時に，回答者の数だけそのプロットもされています。先ほどの3変数でも，調査データの場合は100人，200人とデータを集めることがあるでしょうが，変数の空間の中に投影された100〜200の回答者の点を考えると，もう何が何やらわからない，ということになるかもしれません。

しかし大事なのは，リッカート法で何を測っているか，何が見たかったのか，です。便利だからといって理論的根拠もなく中途半端な方法で，お決まりの分析手法しか使えないようになってしまっては，本質を欠いているのではないでしょうか。
双対尺度法は，カテゴリに対して分析者が主体的に「意味づけする」ような数字を与えます。順序データでも一対比較砲のデータであっても構いません。また強制分類法のように判断基準も能動的に与えることもできます。

紙とペンで得られた反応の分析も，様々な可能性が考えられるはずなのです。
\end{document}